\section{系统地图、Q\&A 限制与研究议程}

四个案例完成后，课程把它们重新映射到一台机器人系统，并在问答中主动暴露上限：skill library 仍小，imitation 只是起点，compositional scaling 没有阈值，六帧视觉上下文又已经接近模型容量。

\subsection{What's Next：先回到系统层}

Agenda 最后一项不是再介绍一篇论文，而是把 local successes 组合起来。任何“机器人 foundation model”声明都应回答：skill learning、planning、low-level control、data augmentation 和 object representation 分别由谁负责，接口如何传递，失败如何反馈。

\lecturefigure{slide-42-agenda-whats-next.jpg}{Agenda：从四篇工作返回完整机器人系统}{官方录像恢复幻灯片 V042；Stanford CS25 Lecture 15}

\subsection{Component Map：不同 Foundation Models 放在哪里}

系统表把 RT-1 放在 skill learning，SayCan / Inner Monologue 放在 planning，Code as Policies 放在 low-level control，DIAL 放在 data augmentation，NLMap 放在 object-centric representations。顶端 recipe 强调这些方法共同依赖 large diverse offline data、high-capacity architecture 和 language glue。

\lecturefigure{slide-43-system-component-map.jpg}{系统 component map：RT-1、SayCan、Inner Monologue、Code as Policies、DIAL 与 NLMap}{官方录像恢复幻灯片 V043；Stanford CS25 Lecture 15}

\begin{importantbox}{Language 是 API，不是单体智能}
语言让 skill、plan、scene description 和 label 共享可读接口，因此外部模型容易接入；但每个模块仍需要独立 data、evaluation 与 failure handling。把所有模块合成一个 decoder 可能简化接口，却不会自动消除 grounding 与实时控制问题。
\end{importantbox}

\subsection{2016--2023：问题如何变化}

timeline 把三阶段问题写在横轴下：2016--2021 问端到端真实机器人如何学习；2021--2022 问如何扩展到许多 tasks；2022--2023 问如何借 foundation models 加速 robotics。右侧列出的 SayCan、Inner Monologue、DIAL、Code as Policies 与 NLMap 是系统部件，不是一个统一 end-to-end model。

\lecturefigure{slide-44-robotics-progress-timeline.jpg}{机器人研究进展：从 end-to-end、multi-task scaling 到 foundation-model leverage}{官方录像恢复幻灯片 V044；Stanford CS25 Lecture 15}

\subsection{Closing 与公开项目}

时间线已经把 RT-1、SayCan、Inner Monologue 与 DIAL 放回各自解决的问题，结尾则把这些系统重新落到可复查的论文、代码和演示上。本节不是再增加一种新方法，而是检查整堂课留下了哪些可继续验证的入口。看最后一页时应把公开项目与前文的 data、planning、feedback 和 grounding 模块一一对应，同时避免把精心挑选的 demo 当成长期自主运行的完整证明。

Thank-you slide 保留四个项目入口和一个清理可乐的机器人示例。它提醒读者本讲的证据来自可查论文、代码与 demo，而不是抽象愿景。演示中的长指令仍被分解为有限技能序列，体现了 language planning 与 skill library 的边界。

\lecturefigure{slide-45-thank-you.jpg}{结语：SayCan、Inner Monologue、RT-1、DIAL 项目与清理任务演示}{官方录像恢复幻灯片 V045；Stanford CS25 Lecture 15}

\subsection{Q\&A 一：Generalization 没有统一阈值}

学生询问组合泛化需要多大规模。讲者回答取决于 task definition、dataset scale 和 concept，当前没有一个数字能预测跨颜色、物体和技能的 emergence；他个人猜测可能还差一到两个数量级的数据，但明确把这作为主观判断。

\begin{warningbox}{不要把“一两个数量级”写成 Scaling Law}
这句话来自问答中的 rough intuition，没有控制实验、横轴定义或置信区间。它能表达“当前规模仍小”，不能用来预测某个参数量或 episode count 必然出现通用机器人能力。
\end{warningbox}

\subsection{Q\&A 二：SayCan 的瓶颈是 Grounded Skills}

SayCan 的 affordance/value functions 只能给有限 skill library 打分。如果机器人只有三项技能，planner 最多产生三项技能的组合；扩大语言模型不会增加新的动作 primitive。高层 instruction coverage 因而取决于低层技能的数量、精度和状态覆盖。

\teachervoice{讲者称这是 SayCan 的“100\% bottleneck”：planner 的 buffet 里有多少菜，决定它能组合多少目标。这个比喻比“LLM 负责 planning”更准确，因为它把能力上限放回机器人数据与 controller。}

\subsection{Q\&A 三：Imitation 是 Existence Proof，不是终局}

RT-1 当时主要使用 imitation loss。讲者作为 RL 研究者希望加入 offline improvement 或 RL，但选择 multi-task imitation 是因为它稳定、可扩展、已经证明大规模真实数据能工作。工程上先建立可靠 base policy，再研究更复杂 optimization，比一开始同时解决所有问题更可验证。

\begin{knowledgebox}{为什么不直接 Online RL}
真实机器人 reward 稀疏、失败昂贵、数据并行受硬件限制。Offline imitation 牺牲主动探索，换来稳定训练和可复用 dataset。未来 RL 若要加入，仍需处理 distribution shift、安全约束与离线 value extrapolation。
\end{knowledgebox}

\subsection{Q\&A 四：六帧已经接近 Context Limit}

长任务需要分钟级记忆，但高维图像 token 很快耗尽 context。RT-1 的六帧只覆盖短期运动，无法存储整条轨迹或 few-shot demonstration。讲者提出 retrieval transformer 或其他 memory mechanism 作为可能方向，并明确说当前把完整轨迹塞入 context 并不现实。

\teachervoice{最后的限制非常具体：不是“模型还不够聪明”，而是六张图像就已接近 context budget。长时具身任务需要压缩、检索、状态摘要或外部 memory；简单扩大 prompt 会同时推高 latency 和实时控制风险。}

\subsection{本章小结}

系统地图显示四篇工作分别填补 skill、plan、feedback 和 data 接口；问答则指出它们没有自动合成通用机器人。真正的剩余瓶颈是 grounded skill coverage、学习算法、可预测的 generalization 和长时 memory。

